% from ega4-2-fr.tex
\oldpage[IV-2]{196}

La propriété $\boldsymbol P$ étant définie comme dans (7.3.1), nous dirons
que \emph{les fibres formelles de $A$ ont la propriété $\boldsymbol P$},
ou que $A$ est un \emph{$\boldsymbol P$-anneau}, si, pour tout
$x\in\operatorname{Spec}(A)$,
$\boldsymbol P(f^{-1}(x),k(x))$ est vraie. Comme $f$ est \emph{plat}, il
revient au même de dire que $f$ est un $\boldsymbol P$-morphisme.

\begin{proposition}[7.3.14]
\label{IV.2.s7.3.14-fr}
Soient $A$ un anneau semi-local noethérien,
$\mathfrak m_i$ $(1\leqslant i\leqslant r)$ ses idéaux maximaux, et posons
$A_i=A_{\mathfrak m_i}$; pour que $A$ soit un $\boldsymbol P$-anneau, il
faut et il suffit que chacun des $A_i$ le soit.
\end{proposition}

\begin{proof}
En effet, $\widehat A$ est composé direct des $\widehat A_i$, donc la
fibre formelle de $A$ en un point $x\in\operatorname{Spec}(A)$ est la
\emph{somme} des fibres formelles en $x$ de ceux des $A_i$ tels que
$x\in\operatorname{Spec}(A_i)$.
\end{proof}

\begin{proposition}[7.3.15]
\label{IV.2.s7.3.15-fr}
Soit $A$ un anneau semi-local noethérien.
\begin{enumerate}
\item[\emph{i)}] Si $A$ est un $\boldsymbol P$-anneau, tout anneau quotient de $A$
est un $\boldsymbol P$-anneau.
\item[\emph{ii)}] Si de plus $\boldsymbol P$ satisfait à la condition
$(\mathrm P_{\mathrm{IV}})$ (7.3.6), alors toute $A$-algèbre finie est un
$\boldsymbol P$-anneau.
\end{enumerate}
\end{proposition}

\begin{proof}
Pour tout idéal $\mathfrak a$ de $A$,
$\widehat A\otimes_A(A/\mathfrak a)$ est le complété de
$A/\mathfrak a$, et les fibres formelles de $A/\mathfrak a$ sont donc les
fibres formelles de $A$ aux points de $V(\mathfrak a)$, d'où (i). D'autre
part, si $B$ est une $A$-algèbre finie, donc un anneau semi-local, on a
$\widehat B=B\otimes_A\widehat A$, et (ii) résulte donc de (7.3.7).
\end{proof}

\begin{proposition}[7.3.16]
\label{IV.2.s7.3.16-fr}
On suppose que la propriété $\boldsymbol P$ est du premier (resp. du
second) type, définie à partir d'une propriété $\boldsymbol R$ (7.3.10).
Lorsque $\boldsymbol P$ est du second type, on suppose en outre que la
relation
\begin{quote}
« pour toute extension finie $k'$ de $k$, et tout
$z'\in Z\otimes_k k'$, $\boldsymbol R(\mathscr O_{z'})$ est vraie »
\end{quote}
entraîne $\boldsymbol P(Z,k)$ (ce qui est vérifié pour les exemples
(iv$'$) à (vii$'$) de (7.3.8), en vertu de (6.7.7) et (4.6.1)).

Soit $A$ un anneau semi-local noethérien. Si la propriété
$\boldsymbol R$ vérifie les conditions $(\mathrm R_{\mathrm I})$ et
$(\mathrm R_{\mathrm{II}})$, les propriétés suivantes sont équivalentes :
\begin{enumerate}
\item[\emph{a)}] $A$ est un $\boldsymbol P$-anneau.
\item[\emph{b)}] Pour tout anneau quotient intègre $B$ de $A$ (resp. toute
$A$-algèbre finie intègre $B$) et tout idéal premier $\mathfrak q$ de
$\widehat B$ dont l'image réciproque dans $B$ est 0,
$\boldsymbol R(\widehat B_{\mathfrak q})$ est vraie.
\end{enumerate}
Si de plus $\boldsymbol R$ satisfait à la condition
$(\mathrm R'_{\mathrm I})$, les propriétés a) et b) sont aussi
équivalentes à :
\begin{enumerate}
\item[\emph{c)}] Pour tout anneau quotient $B$ de $A$ (resp. toute $A$-algèbre finie
$B$) si l'on pose $Y=\operatorname{Spec}(B)$,
$Y'=\operatorname{Spec}(\widehat B)$, et si $g:Y'\to Y$ est le morphisme
canonique, on a (avec les notations de (7.3.12))
\begin{equation}
U_{\boldsymbol R}(Y')=g^{-1}(U_{\boldsymbol R}(Y)).
\tag{7.3.16.1}
\label{IV.2.eq.7.3.16.1-fr}
\end{equation}
\end{enumerate}
\end{proposition}

\begin{proof}
L'équivalence de a) et b) est immédiate lorsque $\boldsymbol P$ est du
premier type : en effet, pour tout idéal premier $\mathfrak p$ de $A$, on
a vu que la fibre formelle de $B=A/\mathfrak p$ au point générique (0) de
$\operatorname{Spec}(B)$ n'est autre que la fibre formelle de $A$ au
point $\mathfrak p\in\operatorname{Spec}(A)$ (7.3.13).

Lorsque $\boldsymbol P$ est du second type, l'équivalence de a) et b)
résulte du lemme plus général suivant :
\end{proof}

\begin{lemma}[7.3.16.2]
\label{IV.2.s7.3.16.2-fr}
Soient $A$, $A'$ deux anneaux, $\varphi:A\to A'$ un homomorphisme,
$f:\operatorname{Spec}(A')\to\operatorname{Spec}(A)$ le morphisme
correspondant. Afin que pour tout $x\in\operatorname{Spec}(A)$, toute

